\section{Introduction}
\label{sec:intro}

A Koblitz curve $E_a:y^2+xy=x^3+ax^2+1$ is defined over $\FF_2$ and used
over $\FF_{2^n}$, so the $2$-power Frobenius $\tau(x,y)=(x^2,y^2)$ is an
endomorphism that acts on the prime-order subgroup $G$ as
multiplication by a known scalar $\lambda$. Together with negation it
generates a group $\Gamma$ of order $2n$, and every algorithm for the
discrete logarithm in $G$ can work on $\Gamma$-orbits instead of
points. For Pollard rho that is worth $\sqrt{2n}$
\citep{wienerzuccherato1998,glv2001}; for index calculus with a
$\Gamma$-invariant factor base it is worth $2n$ in the number of
relations and $(2n)^2$ in the linear algebra
\citep{ggmp2020}. Both facts are classical and neither is claimed
here.

What is claimed is narrower and, we think, more useful to someone who
has to build the thing. The $2n$ saving is a statement about counts.
To collect it, a pipeline has to map every decomposed point $P$ to its
orbit representative $\hat P$ \emph{and} to the sign $\epsilon$ and
phase $k$ with $P=\epsilon\tau^{-k}\hat P$, because the relation
coefficient is $\epsilon\lambda^{-k}$ and not $1$. It has to do this
for every candidate the decomposition produces, not only the ones that
turn into relations, since the candidates are what the tables are
probed with. The tables have to be keyed by representatives, so that
they shrink by $2n$ too, and a hit in such a table has to give back
the actual summands, which means recovering a phase that was never
stored. And at the end every logarithm has to verify, which it does
only if nothing was lost on the way. Each of these steps has a cost,
and in the implementations we measured that cost has at various times
been the majority of the pipeline, the majority of the \emph{reference}
the pipeline was compared with, and the difference between a reported
win and a reported loss.

\paragraph{Contributions.}
\begin{enumerate}[leftmargin=1.6em,itemsep=0.25em]
\item \emph{A representation with proofs
  (Sections~\ref{sec:prelim}--\ref{sec:representation}).} The
  signed-Frobenius form $\can(P)=(\hat P,\epsilon,k)$; a proof that
  $\Gamma$ acts freely on $G\setminus\{O\}$ whenever
  $r\nmid\#E_a(\FF_{2^d})$ for proper divisors $d\mid n$, so that for
  prime $n$ the only exceptional words a canonicalizer meets are
  $x\in\FF_2$; a Frobenius-invariant sign chooser $\Tr(y/x)$ for odd
  $n$; and the reconstruction theorem that makes phase-aware relation
  collection correct.
\item \emph{An exact canonicalization shortcut
  (Section~\ref{sec:canonicalization}).} It rests on the standard
  fact that the least rotation of a binary necklace begins at a
  longest cyclic zero run, used here as a candidate filter. Two
  implementations (bit-parallel for one-word fields, set-bit scan for
  two-limb fields) are proved bit-identical to the exhaustive scan,
  including on the exceptional words, and measured at $4.78$ rotations
  per $83$-bit word against $82$.
\item \emph{Phase-aware layouts (Section~\ref{sec:layout}).} The
  $18$-byte root record with an exact $83$-bit key and a $7$-bit phase;
  a hashed-bucket directory with singleton fingerprints; the orbit tag
  that makes recovery from a folded pair table $O(n)$; the row filter
  that stores each orbit of sums once; and the observation that
  canonical keys are skewed toward small values and must be hashed
  before they are indexed, a mistake two independent implementations
  made.
\item \emph{Benchmarks that say what class they are
  (Sections~\ref{sec:pipeline}--\ref{sec:benchmarks}).} Every number
  is footnoted to a committed artefact, read in one unit against a
  matched signed-Frobenius rho and a generic floor, labelled
  engineering, relabelling or accounting, and marked as isolated or
  exploratory. The end-to-end rows all have verified logarithms. None
  is an advance over rho.
\end{enumerate}

\paragraph{What the reader should take away.}
Three things. The normal-basis least rotation, not the squaring
chain, is the right canonicalization, by $8$--$15\times$ per key and
$4\times$ on the whole rho process; and because a signed-Frobenius rho
spends half its instructions canonicalizing, any comparison against
rho that does not match the canonicalization is measuring the wrong
thing (Section~\ref{sec:bench-rho}). The counts a fold divides by $2n$
are not the counts the cost is proportional to: at $n=61$ the folded
pair table stores $121\times$ fewer pairs and costs $1.2\%$ more
operations in the run that made it the default, and in a separate
four-width sweep it is $2.7\%$ dearer than the full tier at one width
and $12\%$ cheaper at another
(Section~\ref{sec:bench-tables}). And the one place a phase-aware base
makes a complete index calculus faster with verified answers,
$6.85\times$ at $n=19$, is a toy whose base does not exist at
$n=131$ (Sections~\ref{sec:bench-e2e},~\ref{sec:related}).

\paragraph{Organisation.}
Section~\ref{sec:prelim} fixes notation and proves the free-action
lemma. Section~\ref{sec:representation} defines the representation and
proves transport and reconstruction. Section~\ref{sec:canonicalization}
gives the zero-gap lemma, the algorithms and their correctness.
Section~\ref{sec:layout} describes the layouts.
Section~\ref{sec:pipeline} fixes where the phase lives in each
pipeline. Section~\ref{sec:benchmarks} reports the measurements.
Section~\ref{sec:related} says what is prior art, and
Section~\ref{sec:limitations} what is not claimed.
Appendix~\ref{app:repro} lists the artefacts.
